\section{Retos futuros}

\subsection{Combinación flexible de interpretabilidad y auditoría}

A medida que los modelos se vuelven más poderosos, la simple replay evaluation ya no basta para garantizar la seguridad; es necesario introducir mechanistic interpretability y audit trails.

\begin{itemize}
    \item Registrar la información del Q gradient que dispara cada action, para facilitar el reverse engineer.
    \item Combinar el audit log con la governance checklist, formando un ciclo de "explain -> sign off".
    \item Al actualizar la policy, activar automated sanity checks automatizados y conservar el diff.
\end{itemize}

\begin{warningbox}{Ignorar la interpretabilidad provoca fallas de cumplimiento}
Si no hay manera de explicar por qué la policy eligió cierta acción, el equipo de gobernanza la marcará como "opaca", lo que retrasará el despliegue. La interpretabilidad y el compliance son dos dimensiones que se sostienen mutuamente.
\end{warningbox}

\subsection{Adaptarse a cambios de regulación y de política}

El docente enfatiza que los cambios de política tienen un impacto enorme en los proyectos de offline RL: restricciones al compartir datos, requisitos de logging e incluso la prohibición de ciertas formas de reward shaping. Es necesario dar seguimiento en tiempo real a la regulación y archivar los cambios documentados.

\subsection{Resumen de la sección}

Los proyectos futuros de offline RL deben vincular la interpretability con la governance, y mantener sincronizada la documentación de política con la documentación técnica, para avanzar con solidez bajo presión regulatoria.

